\documentclass[11pt]{article}
\usepackage[margin=2.4cm]{geometry}
\usepackage[T1]{fontenc}
\usepackage[brazilian,shorthands=off]{babel}
\usepackage{amsmath,amssymb}
\usepackage{booktabs}
\usepackage{hyperref}
\usepackage{url}

\title{Experimento: RNN em VALE3 (estudo)}
\author{Italo Brito Brandão}
\date{18/10/2026}

\begin{document}
\maketitle

\noindent
Isto é um caderno de estudo. Não é recomendação de compra ou venda.

\section{Pergunta}

Dada uma janela dos últimos \(N\) pregões de VALE3, a rede deve dizer se o \emph{fechamento do dia seguinte} sobe ou desce.

\section{Dados}

Fonte: API pública \href{https://brapi.dev}{brapi.dev}, ticker \texttt{VALE3}, intervalo diário, \texttt{range=1y}.

Arquivos nesta pasta:

\begin{center}
\begin{tabular}{ll}
  \toprule
  arquivo & papel \\
  \midrule
  \texttt{fetch\_vale3.py} & baixa a série e grava o CSV \\
  \texttt{vale3\_1y.csv} & 247 pregões, 20/10/2025 a 16/10/2026 \\
  \texttt{experimento.tex} & este texto \\
  \bottomrule
\end{tabular}
\end{center}

Para atualizar:

\begin{verbatim}
python3 vale3/fetch_vale3.py
\end{verbatim}

\subsection{Colunas}

\begin{center}
\begin{tabular}{ll}
  \toprule
  coluna & significado \\
  \midrule
  \texttt{date} & data do pregão (UTC convertida) \\
  \texttt{open, high, low, close} & preços do dia \\
  \texttt{adj\_close} & fechamento ajustado \\
  \texttt{volume} & negócios \\
  \texttt{next\_close} & fechamento do pregão seguinte \\
  \texttt{label} & \(1\) se \texttt{next\_close} \(>\) \texttt{close}, senão \(0\) \\
  \bottomrule
\end{tabular}
\end{center}

A última linha não tem \texttt{next\_close} nem \texttt{label}: não existe o dia seguinte na amostra. Restam 246 exemplos rotulados: 130 altas, 116 baixas. Fechamento no primeiro dia \(60{,}90\); no último \(66{,}21\).

\section{Rótulo}

No pregão \(t\),

\begin{equation}
  y_t =
  \begin{cases}
    1 & \text{se } C_{t+1} > C_t \quad \text{(sobe)} \\
    0 & \text{caso contrário} \quad \text{(não sobe).}
  \end{cases}
\end{equation}

\(C_t\) é o \texttt{close} do dia \(t\). Isso é classificação de um passo, não um backtest com corretagem, slip ou posição.

\section{Entrada da RNN}

Cada exemplo é uma sequência de \(N\) vetores, um por dia. Um vetor mínimo:

\begin{equation}
  x_t =
  \bigl(
    \log C_t - \log C_{t-1},\;
    \tfrac{H_t - L_t}{C_t},\;
    \log(1+V_t)
  \bigr)
\end{equation}

Preço cru não entra: a rede memorizaria o nível (R\$~60 vs R\$~66), não o movimento. \(N\) pequeno (5 a 20) cabe em 246 pontos; \(N=60\) deixa poucos trechos.

A RNN lê \(x_{t-N+1},\ldots,x_t\) e produz \(h_t\) com a mesma recorrência do caderno \texttt{RNN.tex}:

\begin{equation}
  h_k = \tanh(W_{xh} x_k + W_{hh} h_{k-1} + b_h),
  \qquad
  \hat{y}_t = \sigma(W_{hy} h_t + b_y).
\end{equation}

Só o último estado vira probabilidade de alta. A perda é entropia cruzada binária entre \(\hat{y}_t\) e \(y_t\).

\section{Corte treino / teste}

A série é temporal: não embaralhar dias. Um corte simples é os primeiros 80\% dos exemplos rotulados para treino (ordem cronológica) e o resto para teste. Qualquer vazamento de \(C_{t+1}\) para dentro de \(x_t\) invalida o número.

\section{O que este experimento \emph{não} mostra}

\begin{itemize}
  \item 246 pontos é pouco para uma RNN. Um baseline (sempre prever a classe majoritária, ou ``sobe se ontem subiu'') precisa aparecer na mesma tabela.
  \item Mercado muda de regime. Acurácia no último trimestre de 2026 não se transfere para 2027.
  \item Sem custo de operação, ``comprar todo dia que a rede diz 1'' não é PnL.
  \item brapi agrega fonte pública. Para estudo serve; não é feed da B3 em tempo real.
\end{itemize}

\section{Próximo passo de implementação}

\begin{enumerate}
  \item Carregar \texttt{vale3\_1y.csv}, montar janelas de tamanho \(N\), descartar a última linha.
  \item Treinar uma RNN (ou GRU) de estado pequeno, um passo de classificação no fim.
  \item Reportar acurácia e matriz \(2\times 2\) no bloco de teste, ao lado do baseline.
\end{enumerate}

\end{document}
